\section{Data}
\label{sec:data}

Our dataset comprises 30-minute intraday observations of the S\&P~500
beginning in January 2005. The forecasting target is realized
\emph{variance}, measured as the sum of squared 5-minute log-returns within
each 30-minute bar ($RV_t = \texttt{sumret2}$; see
Equation~\eqref{eq:rv}). We use ``realized variance'' throughout for the
target itself and reserve ``volatility'' for its square root, on which the
models are estimated (Section~\ref{sec:semantic_transforms}).

\subsection{Exogenous Variables}
\label{sec:exog_vars}

In addition to the target series, we incorporate variables spanning several
information categories:

\begin{itemize}
    \item \textbf{Higher-order return moments}: bipower variation, realized semi-variance, autocovariance, absolute returns, cubic returns, and quartic returns.
    \item \textbf{Liquidity and microstructure}: trading volume, turnover (total, buy-initiated, sell-initiated), effective spread, quoted spread, and number of trades.
    \item \textbf{Cross-sectional stock factors}: equal-weighted and value-weighted aggregates of stock-level return moments, bipower variation, semi-variance, and absolute returns.
    \item \textbf{S\&P 500 and sentiment}: SPY trading activity (turnover, buy/sell turnover) and StockTwits social media metrics (attention, sentiment score, sentiment count).
    \item \textbf{Implied volatility}: VIX, VVIX, and VIX3M.
    \item \textbf{Volatility demand}: Options order flow indicators (SPX open/close, all options open/close, open only).
\end{itemize}

\subsection{Feature Accounting}
\label{sec:feature_accounting}

Because several claims in this paper are per-column
(Section~\ref{sec:fwl}) and because the design width is itself part of the
dense-but-weak argument, the mapping from raw series to design columns must be
explicit rather than implied. Three counts circulate in this project---41
exogenous channels, a 359-column exogenous block, and a 529-column design
matrix---and they cannot all be reconciled by the stated expansion rule of
$|\mathcal{L}| = 6$ rolling means per channel, since $41 \times 6 = 246$.
\pending{a per-bucket table giving channels, the lag set applied to each,
derived columns, and the reconciliation to 529; no per-column claim in this
paper should be read as final until this table exists, and the ratio in
Equation~\eqref{eq:per_feature} depends on it directly}

\subsection{Data Cleaning}
\label{sec:cleaning}

The raw data is gridded to a regular 30-minute frequency and filtered to
exclude non-trading periods: Friday evenings after 20:00, all of Saturday, and
Sunday mornings before 18:30. Overnight gaps in stock factors and volatility
demand measures are forward-filled with neutral values. After cleaning and lag
feature generation---which requires a burn-in of 3{,}125 bars for the longest
lag on the base-5 ladder---the frozen evaluation panel comprises exactly
$218{,}934$ observations, and every arm reported in this paper is scored on
those identical rows.

\subsection{A Defect in That Fill}

The sentence above---``forward-filled with neutral values''---understates what
the pipeline did, and the understatement was expensive enough to report as a
finding rather than a correction. Section~\ref{sec:stale_data} is that report.
It is placed immediately after this section because no result in the rest of
the paper is interpretable without it: the panel described above is not the
panel most of our earlier drafts were fitted on.
